红外反射光谱中的周期性干涉条纹能够反映外延层的光学厚度。题目给出SiC晶圆和Si晶圆在$10^\circ$、$15^\circ$入射角下的波数--反射率数据，要求从分层介质中的界面反射与层内传播出发建立测厚模型，确定两类晶圆的外延层厚度，并分析多次内部反射对测量结果的影响\cite{problem}。

问题一要求在空气--外延层--衬底三层结构中，仅保留表面反射光和第一束内部反射光，建立厚度、波数、入射角以及材料折射率与反射率之间的双光束干涉关系，为实测光谱反演提供理论模型。

问题二要求利用附件1、2给出的SiC双角度反射光谱，在双光束模型基础上反演外延层厚度，并结合两角度拟合结果及不同局部极小解比较厚度结果的稳定性，给出SiC晶圆的综合厚度。

问题三要求推导多次内部反射形成稳定多光束干涉的条件，利用附件3、4的Si双角度光谱建立多光束反演模型并确定硅外延层厚度；同时量化高阶反射对反射光谱及厚度结果的影响，并据此判断SiC厚度是否需要进一步考虑多光束修正。
